%% ---------------------------------------------------------------------
\section{非微扰地位与可解模型}\label{sec:power-np}

前几节已经多次碰到同一个问题：领头幂次之外的一切，最终都依赖于我们对 QCD 的非微扰结构知道多少。本节先把全书的非微扰地位集中说清楚，然后介绍三个可以在其中真正非微扰地检验 OPE 与部分子图像的“实验室”。

\begin{remark}[非微扰地位]\label{rem:power-np-status}
\begin{enumerate}[label=(\roman*),leftmargin=1.8em,itemsep=2pt,topsep=2pt]
\item 本书的每一个因子化定理，要么是重整化微扰论中关于红外正规化的部分子态的陈述（\tagR，$R_{\mathrm{PT}}$，或 \tagA），要么是关于微扰论加上编号假设 \cref{asm:intro-wightman}、\cref{asm:intro-composite}、\cref{asm:intro-np-ope}、\cref{asm:intro-af-coeff}、\cref{asm:intro-growth}、\cref{asm:intro-bridge} 的陈述（\cref{def:intro-MA}）。QCD 尚未作为四维量子场论被构造（四维 Yang--Mills 理论的存在性与质量间隙是千禧年问题之一）；因此，“对强子严格”永远意味着“在所列的非微扰假设之下证明完整”（$R_{\mathrm{cond}}$），而 [R] 本身从不意味着对真实世界 QCD 的数学证明。% M28
\item 在 \cref{asm:intro-wightman} 之下有一批真正的非微扰定理：谱函数的解析性与正定性（\cref{prop:ee-analytic}），由 Wightman 公理得到的有限减除色散关系（\cref{prop:dis-dispersion-np}），以及在此基础上、另加非微扰 OPE 的 DIS 矩定理（\cref{thm:dis-moments-np}）；部分子分布的支撑（\cref{prop:pdf-support-hadron}，\tagA）与求和规则（\cref{prop:pdf-sum-rules}）也属于这一层次。最干净的非微扰连接是格点上定域扭度 2 矩阵元的计算（第~22~章的 \cref{prop:lattice-moments}）：它提供矩阵元，不检验强子截面的因子化。
\item 幂次压低的逐阶含义：在微扰论的每一阶，余项在积分范数下被 $O(\lambda^aL^K)$ 控制，并对维数正规子一致（\cref{def:method-power-suppressed}），其中常数随阶数增长；这些逐阶的界不能对 $\as$ 求和。展开是渐近的而非收敛的，而重整化子（\cref{sec:power-renormalons}）表明，在领头歧义之外，微扰级数与幂次修正不能被干净地分开。
\item 一般方法给出的指数可以离真实的幂次很远。在单圈离壳三角图中，一般的误差估计（\cref{lem:appN-error-table}）在软区只保证相对 $O(\lambda^{1/9})$ 的误差，而精确的余项是完整的一个幂次、并且在次领头幂次没有对数：$\int r=-\frac{4m^2}{Q^4}+\frac{m^4}{Q^6}\big(L^2-15+\frac{\pi^2}{3}\big)+O\big(\frac{m^6}{Q^8}L^2\big)$，$L=\ln(Q^2/m^2)$（\cref{prop:appN-triangle-exact}(d)）；在纤维化积分范数下，余项被 $C\lambda^2L^2$ 控制（\cref{prop:appN-triangle-cs}）。
\item 在公理化场论中，OPE 可以在一个相空间条件下被非微扰地导出\cite{Bostelmann:2005rr}（据该文摘要）；这一条件对 QCD 尚未建立。
\item 非微扰效应可以以 $\as$ 的任何幂次都看不见的形式出现：瞬子贡献 $\propto\ee^{-2\pi/\as(Q)}=(\Lambda/Q)^{\beta_0}$（单圈），对偶性破坏在 OPE 的幂次展开中不可见（\cref{conj:ee-duality}）；二者与 \cref{asm:intro-af-coeff} 作为渐近陈述是相容的。
\item 可解模型给出两种相反的教训：二维 $U(N)$ Gross--Neveu 模型中，扭度展开在 $1/N$ 的领头阶绝对收敛；而超可重整的二维 $O(N)$ 模型中，动量空间的幂次展开因阶乘增强而发散，坐标空间的展开则收敛（第~23~章的 \cref{thm:yzliu-gn-convergence}、\cref{prop:yzliu-on-factorial}）。它们是模型结果，不证明 QCD 中扭度展开的收敛或发散。
\end{enumerate}
\end{remark}

\subsection{微扰 \texorpdfstring{$\varphi^4_4$}{phi4} 中 OPE 的收敛}

OPE 通常被理解为短距离的渐近展开。在一个具体的模型中，可以证明它在每一阶微扰论上实际上收敛，而且在任意有限距离上收敛。

\begin{theorem}[微扰 $\varphi^4_4$ 中 OPE 的收敛性 \tagR]\label{thm:power-ope-convergence}
（$R_{\mathrm{PT}}$：在重整化微扰论中对红外正规化的部分子态严格。）（本定理只涉及 $\varphi^4_4$ 模型的微扰论。）考虑四维欧氏空间中质量 $m>0$ 的厄米标量场 $g\varphi^4$ 理论，取定圈数 $l$。复合场 $O_A=\partial^{w_1}\varphi\cdots\partial^{w_n}\varphi$ 的量纲为 $[A]=n+\sum_i|w_i|$；旁观场取为 $\varphi(f_{p_i})=\int\dd^4x\,\varphi(x)f_{p_i}(x)$，其中 $f_{p_i}$ 光滑，$\hat f_{p_i}$ 的支撑在以 $p_i$ 为心、半径 $\epsilon$ 的球内。把 OPE $O_A(x)O_B(0)\sim\sum_C\mathcal C^C_{AB}(x)O_C(0)$ 中的求和限于 $[C]-[A]-[B]\le\Delta$，则在 $l$ 圈阶有
\begin{multline}\label{eq:power-hk}
\Big|\Big\langle O_A(x)O_B(0)\prod_{i=1}^n\varphi(f_{p_i})\Big\rangle-\sum_C\mathcal C^C_{AB}(x)\Big\langle O_C(0)\prod_{i=1}^n\varphi(f_{p_i})\Big\rangle\Big|
\le m^{[A]+[B]+n}\sqrt{[A]!\,[B]!}\,\tilde K^{[A]+[B]}\prod_i\sup|\hat f_{p_i}|\\
\times\Big(\sup\big(1,\tfrac{|\vec p|_n}{m}\big)\Big)^{2([A]+[B])(n+2l+1)+3n}\sum_{\lambda=0}^{n/2+2l}\frac{\ln^\lambda\sup(1,\frac{|\vec p|_n}{m})}{2^\lambda\lambda!}
\cdot\frac{1}{\sqrt{\Delta!}}\Big(\tilde K\,m|x|\sup\big(1,\tfrac{|\vec p|_n}{m}\big)^{n+2l+1}\Big)^{\Delta},
\end{multline}
其中 $\tilde K$ 只依赖于 $n$ 与 $l$，$|\vec p|_n\equiv\sup_{J\subset\{1,\dots,n\}}|\sum_{i\in J}p_i|$\cite[定理~3]{Hollands:2011gf}。由于对任意 $c$ 有 $c^\Delta/\sqrt{\Delta!}\to0$，在每一圈阶上 OPE 在每个固定的 $|x|$ 处收敛，而不只是渐近地成立。
% VF-FRO-4
% uses: none
\end{theorem}

\begin{proofsketch}
证明\cite{Hollands:2011gf}采用 Wilson--Wegner--Polchinski 重整化群流方程：引入红外截断 $\Lambda$ 与紫外截断 $\Lambda_0$，对带复合算符插入的连通截肢 Green 函数推导关于 $\Lambda$ 的流方程，在圈数与外腿数上归纳地给出对 $\Lambda_0$ 一致的界；关键在于一个足够精细、能在归纳中自我复制的归纳假设，它同时控制 OPE 的余项。一个附带结果是有效作用量的“梯度展开”在每一圈阶收敛。完整证明见\cite[定理~3 及其证明]{Hollands:2011gf}；无质量 $\varphi^4_4$ 的同一结论见\cite{Holland:2014pna}（据该文摘要），弯曲时空中微扰 OPE 的构造见\cite{Hollands:2006ag}（据该文摘要）。
\end{proofsketch}

这一定理的意义与限度都需要说清楚。收敛性依赖于旁观场所代表的“态”：$\hat f_{p_i}$ 的支撑若不紧，高能行为越差，收敛性越差；圈数展开本身的收敛性没有涉及；而在闵氏空间中只能期望对类空分离收敛\cite{Hollands:2011gf}。它是一个模型陈述：对含夸克的 QCD，微扰 OPE 是全阶论证（\cref{thm:dis-ope}，\tagA），其收敛性是猜想；更多严格的 OPE 结果（纯 Yang--Mills 中 OPE 在规范不变算符上闭合、共形场论中的光线 OPE 等）见第~24~章（\cref{prop:front-ope-scalar}、\cref{prop:front-ope-ym}、\cref{conj:front-ope-convergence}、\cref{rem:front-ope-qcd-status}）。

\subsection{二维大 \texorpdfstring{$N_c$}{Nc} 的 't Hooft 模型}

二维 QCD 在 $N_c\to\infty$ 极限中是可解的\cite{tHooft1974}，而且它恰好是一个检验“禁闭的束缚态是否仍有部分子图像”的理想场所：夸克被禁闭，介子谱是离散的，但深度非弹散射仍然可以精确计算。

\paragraph{模型与 Feynman 规则。}
取二维 Minkowski 空间的光锥坐标 $x^\pm=(x^0\pm x^1)/\sqrt2$，规范群 SU($N_c$)，夸克质量 $m_a$，$D_\mu=\partial_\mu-\ii gA^a_\mu t^a$。在二维中 $g$ 有质量量纲 1，理论是超可重整的。取光锥规范 $A^+=A_-=0$，则 $G^a_{+-}=-\partial_-A^a_+$（非阿贝尔项因只剩一个分量而消失），规范部分的拉氏量为 $\frac12(\partial_-A^a_+)^2$：$A_+$ 没有传播自由度，其“传播子”为 $\langle A^a_+A^b_+\rangle(k)=\ii\delta^{ab}/(k^+)^2$，与 $k^-$ 无关（在光锥时间 $x^+$ 上是瞬时的），$k^+=0$ 处按主值（或截断 $|k^+|>\lambda$）理解；夸克—胶子顶点为 $\ii g\gamma^+t^a$。取 't Hooft 极限 $N_c\to\infty$，$m_0^2\equiv g^2N_c/(2\pi)$ 固定，只剩平面图；由于没有胶子自耦合，平面图就是夸克的彩虹型自能与夸克—反夸克之间的梯形图。

\paragraph{夸克传播子。}
单胶子自能为 $\Sigma(p)=\frac{N_c}{2}(\ii g)^2\ii\int\frac{\dd^2k}{(2\pi)^2}\frac{\gamma^+S(p-k)\gamma^+}{(k^+)^2}$（$t^at^a=C_F\to N_c/2$）。由 $\gamma^+(\sla q+m)\gamma^+=2q^+\gamma^+$，被积函数只通过 $1/[2q^+q^--m^2+\ii0]$ 依赖于 $q^-$；对称地积分 $\int\dd k^-\,[A-2q^+k^-+\ii0]^{-1}=-\ii\pi/(2|q^+|)$，再用 $\int_{|k^+|>\lambda}\dd k^+\,\mathrm{sgn}(p^+-k^+)/(k^+)^2=2/\lambda-2/p^+$（$p^+>0$），得
\begin{equation}\label{eq:power-thooft-sigma}
\Sigma(p)=-\ii\,\frac{g^2N_c}{4\pi}\gamma^+\Big(\frac1\lambda-\frac{1}{p^+}\Big),\qquad
\text{传播子分母}\;2p^+p^--\big(m^2-m_0^2\big)-m_0^2\frac{|p^+|}{\lambda}+\ii0 .
\end{equation}
由于 $\Sigma$ 只依赖于 $p^+$，把内线换成修饰后的传播子不改变 $k^-$ 积分的结果，所以这一单胶子结果就是彩虹求和的精确解。质量平方被移动 $-m_0^2$；与 $\lambda$ 有关的项在 $\lambda\to0$ 时把夸克的极点推到无穷远——单个夸克没有有限能量的态。

\paragraph{束缚态方程。}
对夸克（味 $a$）—反夸克（味 $b$）、总动量 $P$ 的 Bethe--Salpeter 振幅求梯形和，并定义 $\phi(x)=\int\dd k^-\,\Psi(k^+=xP^+,k^-)$。核与 $k^-$ 无关，故对 $k^-$ 的积分可用留数完成：两个夸克传播子的极点只在 $0<x<1$ 时位于实轴两侧，因此 $\phi$ 的支撑在 $[0,1]$ 内。结果是 't Hooft 方程\cite{tHooft1974}
\begin{equation}\label{eq:power-thooft-eq}
\mu^2\phi(x)=\Big(\frac{m_a^2-m_0^2}{x}+\frac{m_b^2-m_0^2}{1-x}\Big)\phi(x)-m_0^2\,\mathrm{f.p.}\!\!\int_0^1\!\dd y\,\frac{\phi(y)}{(y-x)^2},\qquad\mu^2=2P^+P^-,
\end{equation}
其中自能中与 $\lambda$ 有关的项与梯形核中的同类项相消，留下 Hadamard 有限部分 $\mathrm{f.p.}$。利用 $\mathrm{f.p.}\int_0^1\dd y\,(y-x)^{-2}=-\frac1x-\frac1{1-x}$，\eqref{eq:power-thooft-eq}可等价地写成
\begin{equation}\label{eq:power-thooft-pv}
\mu^2\phi(x)=\Big(\frac{m_a^2}{x}+\frac{m_b^2}{1-x}\Big)\phi(x)-m_0^2\,\mathrm{P}\!\!\int_0^1\!\dd y\,\frac{\phi(y)-\phi(x)}{(y-x)^2},
\end{equation}
对应的二次型为
\begin{equation}\label{eq:power-thooft-form}
\langle\phi,H\phi\rangle=\int_0^1\!\dd x\Big(\frac{m_a^2}{x}+\frac{m_b^2}{1-x}\Big)|\phi(x)|^2+\frac{m_0^2}{2}\int_0^1\!\!\int_0^1\!\dd x\,\dd y\,\frac{|\phi(x)-\phi(y)|^2}{(x-y)^2}\;\ge0 .
\end{equation}

\begin{proposition}['t Hooft 模型 \tagP]\label{prop:power-thooft-model}
（在二维 QCD 的 $N_c\to\infty$ 极限中精确；作为关于 QCD 的陈述只是部分的。）在上述模型中：
\begin{enumerate}[label=(\alph*),leftmargin=2em,itemsep=1pt,topsep=2pt]
\item 介子的质量与光锥波函数由\eqref{eq:power-thooft-eq}的本征值问题给出；谱是离散的、非负的，$m_a=m_b=0$ 时最低态无质量且 $\phi_0=1$；端点行为为 $\phi_n(x)\sim x^{\beta_a}$（$x\to0$），$\pi\beta_a\cot(\pi\beta_a)=1-m_a^2/m_0^2$（$x\to1$ 处同理）；高激发态满足 $\mu_n^2\simeq\pi^2m_0^2\,n$。
\item 介子 $n$ 的夸克与反夸克部分子分布为 $q_n(x)=|\phi_n(x)|^2$、$\bar q_n(x)=|\phi_n(1-x)|^2$，支撑在 $[0,1]$ 内，满足数目求和规则 $\int_0^1q_n=1$ 与动量求和规则 $\int_0^1x\,[q_n(x)+\bar q_n(x)]\,\dd x=1$（胶子不携带动量）。
\item 在 Bjorken 极限下，介子的深度非弹结构函数标度化，并且（对末态的离散共振取平均，即作为 $x$ 的分布弱收敛）等于以 (b) 中部分子分布表示的部分子模型结果；修正是 $m_0^2/Q^2$、$m_a^2/Q^2$ 的幂次，没有对数型的标度破坏\cite{Einhorn1976}。
\end{enumerate}
这是一个在规范理论中可以从第一原理出发、非微扰地算出束缚态部分子分布的例子；该模型中的准部分子分布见第~23~章（\cref{prop:yzliu-2dqcd}）。
（本命题只涉及二维大 $N_c$ 模型。）
\uses{def:pdf-quark}
\end{proposition}

\begin{proof}
(a) \eqref{eq:power-thooft-pv}与\eqref{eq:power-thooft-eq}的等价性由上面的有限部分恒等式给出；\eqref{eq:power-thooft-form}由\eqref{eq:power-thooft-pv}关于 $x$ 与 $y$ 对称化得到，两项都非负，故 $\mu^2\ge0$。$m_a=m_b=0$ 时第一项消失，第二项对 $\phi=$ 常数为零，所以 $\phi_0=1$、$\mu_0^2=0$：手征极限下有一个无质量的“$\pi$ 介子”。端点行为：把 $\phi\simeq x^\beta$ 代入\eqref{eq:power-thooft-eq}并只保留 $x\to0$ 时最奇异的项，利用 $\mathrm{f.p.}\int_0^\infty\dd y\,y^\beta(y-x)^{-2}=-\pi\beta\cot(\pi\beta)\,x^{\beta-1}$（$-1<\beta<1$），得 $(m_a^2-m_0^2)+m_0^2\pi\beta\cot(\pi\beta)=0$。对高度振荡的 $\phi_n\simeq\sqrt2\sin(n\pi x)$，主值核作用于 $\ee^{\ii ky}$ 给出 $\pi|k|\ee^{\ii kx}$，于是 $\mu_n^2\simeq\pi^2m_0^2n$。$H$ 是 $L^2[0,1]$ 上的正自伴算符，其预解式是紧的，故谱离散、本征函数完备\cite{tHooft1974,Einhorn1976}。
(b) 在光锥规范中 Wilson 线等于 1，部分子分布（\cref{def:pdf-quark}在二维中的形式）就是光锥时刻 $x^+=0$ 上的数密度算符的矩阵元。$N_c\to\infty$ 时，光锥量子化下介子的价夸克—反夸克分量的波函数就是 $\phi_n$，更高的 Fock 分量被 $1/N_c$ 压低（'t Hooft 解的这一部分子诠释见\cite{Einhorn1976}），因此找到动量分数为 $x$ 的夸克的几率密度是 $|\phi_n(x)|^2$，反夸克携带 $1-x$。归一化 $\int_0^1|\phi_n|^2=1$ 给出数目求和规则，$x+(1-x)=1$ 给出动量求和规则；支撑来自 (a) 中 $\phi$ 的支撑。
(c) 在 Bjorken 极限下 $q^-\to\infty$（$x$ 固定），被击夸克在与之共轭的、长度为 $O(1/q^-)$ 的光锥时间间隔内传播。理论是超可重整的，耦合 $g^2$ 有质量平方的量纲，每一次瞬时相互作用插入到这段硬传播中都带来 $m_0^2/Q^2$ 量级的压低，因此结构函数等于“手提包”项，即 (b) 中分布的部分子模型表达式，修正为 $m_0^2/Q^2$ 与 $m_a^2/Q^2$ 的幂次；没有对数，因为二维中没有对数发散的重整化。末态由离散的介子共振组成，固定 $Q^2$ 时结构函数是 $x$ 的一系列 $\delta$ 函数，相邻共振在 $x$ 上的间距按 $m_0^2/Q^2$ 趋于零，所以标度化是在分布意义下（对任意连续的检验函数）成立的。对全部强子末态求和确实精确地再现部分子模型，这一点已由 Einhorn 明确验证，他同时验证了 Drell--Yan--West 关系与 Bloom--Gilman 对偶\cite{Einhorn1976}；电流振幅的构造与模型的幺正性见\cite{CallanCooteGross1976}。
（\eqref{eq:power-thooft-sigma}--\eqref{eq:power-thooft-form}中的积分、等价形式、正定性、手征极限的零模、端点指数与谱的渐近行为都用符号与数值计算核对过。）
\end{proof}

\paragraph{地位与缺口}
作为关于 QCD 的陈述，这是 \tagP：模型在二维（没有横向自由度、胶子不传播、超可重整，因而没有对数型的标度破坏、没有 Sudakov 对数，也没有 Glauber 问题）并取 $N_c\to\infty$。证明是平面近似中的场论证明，谱的完备性等数学性质依赖于所引的分析\cite{tHooft1974,Einhorn1976}，(c) 中对末态共振的求和由\cite{Einhorn1976}明确完成。它说明的是：禁闭本身并不破坏部分子模型意义下的因子化，束缚态的部分子分布可以从第一原理算出；它不说明四维 QCD 中的任何东西（\cref{asm:intro-bridge} 的地位不因此改变）。

\paragraph{第三个“实验室”：格点。}
格点 QCD 为非微扰 QCD 提供数值证据（\tagP 意义下的证据，而非证明）：它计算局域算符的矩阵元，例如扭度 2 算符的矩（第~22~章），也可以检验重整化子的预言——plaquette 的高阶微扰级数确实显示出量纲 4 重整化子的主导（\cref{prop:power-renormalon-examples}(g)）。格点不检验强子截面的因子化本身，% critic C12
其从欧氏关联函数到光锥分布的匹配又依赖于自己的因子化定理（第~22~章）。

\begin{table}[htbp]
\centering\small
\caption{本章陈述的地位一览。“性质”一栏说明陈述成立的意义（精确、模型中精确、有限阶、标度陈述 E-S 等）。}
\label{tab:power-status}
\begin{tabular}{@{}p{3.3cm}p{3.4cm}p{1.6cm}p{2.2cm}p{3.5cm}@{}}
\toprule
陈述 & 标签 & 等级 & 性质 & 主要缺口 \\
\midrule
DIS 扭度 4（$x$ 空间） & \cref{prop:power-dis-twist4-x} & \tagA & E-S & 次领头幂次区域分析；重整化子混合 \\
SIDIS 型扭度 3 SSA & \cref{prop:power-twist3-sidis} & \tagP & 有限阶 & 无全阶论证 \\
$pp$ 中扭度 3 SSA & \cref{conj:power-twist3-pp} & \tagC & --- & 带色末态 \\
次软定理 & \cref{prop:power-next-to-soft} & \tagR & 树图，精确 & 圈图层次为 \tagA \\
NLP 端点发散 & \cref{prop:power-nlp-endpoint} & \tagP & 逐个过程 & 无一般定理 \\
OPE 基的重排 & \cref{prop:power-rearrangement} & \tagR & 代数恒等式 & --- \\
大 $\beta_0$ Adler 函数 & \cref{prop:power-adler-large-beta0} & \tagP & 极限中精确 & 朴素非阿贝尔化 \\
重整化子的例子；模型 & \cref{prop:power-renormalon-examples} & \tagP & 近似／模型 & 大 $\beta_0$ 之外无控制 \\
重整化子 $\leftrightarrow$ 高扭度 & \cref{conj:power-renormalon-ht} & \tagC & --- & 无非微扰 OPE \\
$\varphi^4_4$ 中 OPE 收敛 & \cref{thm:power-ope-convergence} & \tagR & 每阶显式界 & 模型；QCD 中为猜想 \\
't Hooft 模型 & \cref{prop:power-thooft-model} & \tagP & 模型中精确 & 二维，$N_c\to\infty$ \\
\bottomrule
\end{tabular}
\end{table}

\paragraph{本章小结。}
领头幂次之外，严格性迅速下降（\cref{tab:power-status}）。可以完整证明的是代数性与模型中的陈述：树图层次的次软定理、OPE 基的重排自由、$\varphi^4_4$ 中 OPE 的收敛，以及大 $\beta_0$ 极限与二维大 $N_c$ 模型中的精确结果。DIS 扭度 4 的 $x$ 空间形式有全阶论证；扭度 3 的单自旋不对称与次领头幂次的因子化只在有限阶与特定过程中得到验证；重整化子与高扭度的一般对应以及强子—强子碰撞中的扭度 3 因子化是猜想。贯穿全章的一点是：一旦超出领头幂次，“微扰系数”与“非微扰矩阵元”的划分本身就有与幂次修正同量级的歧义（\cref{prop:power-rearrangement}，\cref{ex:power-renormalon-chain}），因此“精确到幂次修正”只有在幂次修正项被定义之后才有确切含义——这正是\cref{asm:intro-bridge}(Br3) 所表达的内容。前沿方面，线性重整化子与领头重整化子重求和见第~22~章，可解模型中扭度展开的收敛性见第~23~章，次领头幂次的端点因子化与次领头幂次 TMD 因子化见第~24~章。
